\documentclass[10pt,a4paper]{amsart}
\usepackage[margin=19mm]{geometry}
\usepackage{amsmath,amssymb,mathtools}
\usepackage[scr=rsfs,cal=euler]{mathalfa}
\usepackage{xfrac}
\usepackage[unicode,hidelinks,bookmarks=false]{hyperref}
\newcommand{\ie}{i.e.\ }
\newcommand{\cf}{cf.\ }
\newcommand{\resp}{respectively\ }
\newcommand{\Subsection}{\subsection}
\providecommand{\nobreakdash}{\nobreak\hbox{-}}
\setlength{\emergencystretch}{2em}
\multlinegap0pt
\usepackage{bm}
\usepackage{bbm}
\usepackage{dsfont}
\usepackage{xspace}
\usepackage{cancel}
\usepackage{mathtools}
\usepackage[shortlabels]{enumitem}
\usepackage{amsthm}
\makeatletter
\@ifundefined{thm}{\newtheorem{thm}{Thm}}{}
\@ifundefined{lem}{\newtheorem{lem}[thm]{Lemma}}{}
\makeatother
\makeatletter
\newcommand{\CC}{\mathbb{C}}
\newcommand{\GG}{\mathbb{G}}
\newcommand{\RR}{\mathbb{R}}
\expandafter\ifx\csname psmatrix\endcsname\relax
\newenvironment{psmatrix}
  {\left(\begin{smallmatrix}}
\fi
\expandafter\ifx\csname remarks\endcsname\relax
\newenvironment{remarks}
{\medskip\noindent{\it
Remarks.}\begin{list}{{\rm(\arabic{remarkscounter})}
}{\usecounter{remarkscounter}
\renewcommand{\theremarkscounter}{{\bf(\alph{remarkscounter})}}
\setlength{\labelsep}{\fill} \setlength{\leftmargin}{0pt}
\setlength{\itemindent}{\fill}
\setlength{\labelwidth}{\fill}\setlength{\topsep}{0pt}
\setlength{\listparindent}{0pt}}} {\end{list}}
\fi
\makeatother
\title[Geometrization of summation formulae for quadrics]{Geometrization of summation formulae for quadrics}
\author{Chun-Hsien Hsu}
\date{Source: arXiv:2606.01491v1 (CC BY 4.0)}
\begin{document}
\maketitle

\noindent\emph{Source and licence.} Geometrization of summation formulae for quadrics. Version arXiv:2606.01491v1 (CC BY 4.0), distributed under the Creative Commons Attribution 4.0 International licence, as recorded in the arXiv licence metadata for this exact version and shown on the abstract page https://arxiv.org/abs/2606.01491v1. Full rights notice: licenses/arxiv-2606.01491-CC-BY-4.0.txt.\par\smallskip

\noindent\textit{Editorial scope.} Counted unit: the Lem whose source label is \texttt{lem:easy:arch} in the article (source0.tex lines 1314--1335, source label \texttt{lem:easy:arch}), delivered together with the article's own complete original proof of it (source0.tex lines 1337--1407, one \texttt{proof} environment of 71 lines). No mathematics is written, repaired or re-derived: statement and proof are the article's own text line for line. Required context retained uncounted: none. Every definition, display and label of those lines is retained unchanged. Omitted from this package and not redistributed: 1--1313, 1336--1336, 1408--1875. Nothing inside a retained excerpt is omitted or altered: every retained body is delivered exactly as the article's lines are written, so no omission receipt of any kind accompanies this package. Citations: the retained excerpts carry no citation command at all, so no bibliography entry of the article is transcribed and no reference is redistributed. Reference adaptation: none was needed and none was made; every reference inside the delivered unit resolves inside it, so no unresolved reference or citation is produced. Printed slips kept literal: no printed word, symbol, hypothesis or number of the counted statement or of its proof was altered. Layout and class: the article's own class is \texttt{amsart} with packages including amsthm, amsmath, amscd, amsfonts, amssymb, mathtools, geometry, fancybox, pdflscape, rotating, epic, eepic, bbm, longtable; the wrapper is the lane's pinned \texttt{amsart} editorial wrapper at 10pt on A4, which restates the theorem-like environments the retained text needs and adds the article's own macro definitions that the retained text needs (\texttt{\textbackslash CC}, \texttt{\textbackslash GG}, \texttt{\textbackslash RR}), with extra wrapper packages (bm, bbm, dsfont, xspace, cancel, mathtools, enumitem). The delivered numbering is the unit's own. Assets: no retained span contains a figure environment, a graphics-inclusion command, a PDF-page inclusion, a table or a TikZ picture; no image file is copied into this package, the package declares no asset dependency and no image file is redistributed with it. Licence: version arXiv:2606.01491v1 of this article is distributed under the Creative Commons Attribution 4.0 International licence, as recorded in the arXiv licence metadata for this exact version and shown on the abstract page https://arxiv.org/abs/2606.01491v1. A full rights notice is delivered at licenses/arxiv-2606.01491-CC-BY-4.0.txt. Characters: every retained excerpt is pure ASCII, so the delivered engine, XeLaTeX with its default Latin Modern text font, reports no missing character.
\begin{lem}\label{lem:easy:arch}
\begin{enumerate}[label=(\roman*)]
    \item  Suppose $\ell\ge 3$. For $f\in W_\ell'$ we have
    \begin{align*}
        \tilde{c}_\ell(f)=\begin{cases}
            I(f)(0) &\textrm{if } \ell \textrm{ is odd and }F=\RR,\\
            0 & \textrm{otherwise.}
        \end{cases}
    \end{align*}     
    \item Suppose $\ell=2$. For $f\in\mathcal{S}(V_2(F)\oplus F^2),$ $x\in X_{2}(F)^1\cong \mathbb{P}X_2(F)$ 
    \begin{align*} 
        I(f)(ax)=-\tilde{c}_2(f)\log|a|+\tilde{a}_2(f)(x)+o(1) \quad\quad \textrm{ as } |a|\to 0
    \end{align*}
    for some function $\tilde{a}_2(f)\in C^\infty(\mathbb{P}X_2(F))$, and the difference
    \begin{align*}
 Z_{2}(f,s)-\frac{1}{s}\tilde{c}_2(f)
    \end{align*}
    is holomorphic for $\mathrm{Re}(s)>-1/2$. Both $\tilde{c}_2(f)$ and $\tilde{a}_2(f)$ depend only on $I(f)$.
    \item For $f\in \widetilde{\mathcal{S}}(X_1(F)),$ the integral defining $I(f)(x)$ is absolutely convergent for all $x\in X_1(F)$. We have $\tilde{c}_1(f)=\zeta(1)^{-1}I(f)(0).$
\end{enumerate}

\end{lem}

\begin{proof}
    Let $f\in \mathcal{S}(V_{\ell+1}(F)).$ We may assume $f$ is $r_\ell(K)$-invariant. Therefore, for $x\in X_\ell^\circ(F)$
    \begin{align*}
        I(f)(x)&=\int_{F^\times} |t|^{2-\ell} f(t^{-1}x,0,t)d^\times t
    \end{align*}
    and 
    \begin{align*}
        Z_\ell(f,s)=\int_{F^\times} f(0_{\underline{2\ell}},0,t)|t|^sd^\times t.
    \end{align*}
    Both integral converges absolutely when $\ell=s=1,$ in which case we can take $x=0$. This proves (iii).
    
   Suppose $\ell\ge 3$. Since these two operators are continuous in $f$, to compute both terms we now assume $f(v,0,t)=f_1(v)f_2(t)$ for some $f_1\in \mathcal{S}(F^{2\ell})$ and  $f_2\in \mathcal{S}(F)$.  We will prove the lemma for $F=\RR.$ The case $F=\CC$ is analogous by identifying $\CC=\RR^2,$ so we leave it to the reader. 
   
    Choose $r>0$ and complex numbers $b_0,\ldots, b_{\ell-2}$ such that 
   \begin{align*}
       \bigg|f_2(t)-\sum_{n=0}^{\ell-2}b_nt^n\bigg|\le |t|^{\ell-2+1/4} \textrm{ for all } |t|\le r. 
   \end{align*}
   Then for $a\in F^\times$
   \begin{align*}
       I(f)(ax)&=\int_{|t|>r} |t|^{2-\ell}f_1(t^{-1}ax)f_2(t) d^\times t+\int_{|t|\le r} |t|^{2-\ell}f_1(t^{-1}ax)\left(f_2(t)-\sum_{n=0}^{\ell-2} b_nt^n\right) d^\times t\\
       &+\sum_{n=0}^{\ell-2}\int_{|t|\le r} b_nt^n|t|^{2-\ell} f_1(t^{-1}ax)d^\times t.
   \end{align*}
   The first two terms are absolutely convergent and bounded above by a constant independent of $a$ and $x$. For the last term, by changing variables we obtain
   \begin{align*}
       \sum_{n=0}^{\ell-2} b_n a^n|a|^{2-\ell}\int_{|t|\le |a|^{-1}r}t^n|t|^{2-\ell} f_1(t^{-1}x)d^\times t.
   \end{align*}
    Therefore, if $f\in W_\ell'$ so that $I(f)(0)=\lim_{|a|\to 0}I(f)(ax)$ exists, each summand above is zero except possibly for $n=\ell-2$. In this case, we arrive at
    \begin{align*}
        b_{\ell-2}\mathrm{sgn}^{\ell}(a)\int_{|t|\le |a|^{-1}r} \mathrm{sgn}^\ell(t)f_1(t^{-1} x)d^\times  t.
    \end{align*}
    This term must be zero if $\ell$ is odd. For $\ell$ even, it must be $f_1(0)=0.$ 
    
    When $\ell$ is odd, we have 
    \begin{align*}
        I(f)(0)&=f_1(0)\left(\int_{|t|>r} |t|^{2-\ell}f_2(t) d^\times t+\int_{|t|\le r} |t|^{2-\ell}\left(f_2(t)-b_{\ell-2}t^{\ell-2}\right) d^\times t\right).
    \end{align*}
    On the other hand, for any $\ell\ge 3$ we have
    \begin{align*}
        Z_\ell(f,s)=f_1(0)\left(\int_{|t|>r} f_2(t)|t|^{s}d^\times+\int_{|t|>r} (f_2(t)-b_{\ell-2}t^{\ell-2})|t|^sd^\times t+\int_{|t|<r} b_{\ell-2}t^{\ell-2}|t|^s d^\times t\right).
    \end{align*}
    This is zero when $\ell$ is even. When $\ell$ is odd, the last integral vanishes, so $I(f)(0)=Z_\ell(f,2-\ell).$ This proves (i). 
    
    Suppose $\ell=2$. Let $x\in X_2(F)^1 $ and $|a|<1$. Write
    \begin{align*}
        I(f)(ax)&=\int_{|t|>1} f(t^{-1}ax,0,t)d^\times t+ \int_{|t|\le 1} f(t^{-1}ax,0,t)-f(t^{-1}ax,0,0)d^\times t\\
        &+\int_{|t|\le |a|^{-1}} f(t^{-1}x,0,0)d^\times t.
    \end{align*}
    The last term can be written as
    \begin{align*}
         &\int_{|t|< 1} f(t^{-1}x,0,0)d^\times t+\int_{1\le |t|\le |a|^{-1}} f(t^{-1}x,0,0)d^\times t\\
    &=\int_{|t|> 1} f(tx,0,0)d^\times t+\int_{1\ge |t|\ge|a|} f(tx,0,0)-f(0)d^\times t -\frac{\mathrm{vol}(K_{\GG_m})}{[F:\RR]}f(0)\log |a|.
    \end{align*}
    Observe that
    \begin{align*}
        \lim_{|a|\to 0} I(f)(ax)+\frac{\mathrm{vol}(K_{\GG_m})}{[F:\RR]}f(0)\log |a| 
    \end{align*}
    is well-defined, and we denote this as $\tilde{a}_2(f)(x)$. Explicitly,
    \begin{align}\label{eq:a2:arch}
    \begin{split}
        \tilde{a}_2(f)(x)&=\int_{|t|>1} f(0_{\underline{4}},0,t)d^\times t+ \int_{|t|\le 1} f(0_{\underline{4}},0,t)-f(0)d^\times t\\
        &+\int_{|t|> 1} f(tx,0,0)d^\times t+\int_{|t|\le 1} f(tx,0,0)-f(0)d^\times t.
    \end{split}
    \end{align}
    
    
    On the other hand,
    \begin{align}\label{eq:Z2compute}
        Z_2(f,s)=\int_{|t|\ge 1} f(0_{\underline{4}},0,t)|t|^s d^\times t+\int_{|t|<1} (f(0_{\underline{4}},0,t)-f(0))|t|^s d^\times t+\int_{|t|<1} f(0)|t|^s d^\times t.
    \end{align}
    The first two integrals are absolutely convergent for $\mathrm{Re}(s)>-1/2$. The last term above is $\frac{1}{[F:\RR]s}\mathrm{vol}(K_{\GG_m})f(0),$ so $\tilde{c}_2(I(f))=\frac{\mathrm{vol}(K_{\GG_m})}{[F:\RR]}f(0).$ This justifies (ii).
\end{proof}

\end{document}
